\documentclass{article}
\usepackage[utf8]{inputenc}
\usepackage[T1]{fontenc}
\usepackage{geometry}
\usepackage{graphicx}
\usepackage{float}
\usepackage{hyperref}
\usepackage[english,indonesian]{babel}

\title{Optimasi Pemrosesan Template LaTeX Menggunakan Pendekatan Berbasis AI}
\author{Academic Support Agent}
\date{\today}

\begin{document}
\maketitle

\begin{abstract}

\end{abstract}

\chapter{PENDAHULUAN}


\chapter{TINJAUAN PUSTAKA}


% Original content (commented out):
% \documentclass[10pt,a4paper]{article}
\include{Formato_TFM}
\usepackage[utf8]{inputenc}
\usepackage{enumerate}
\include{Datos_TFM}
%------------------------------------------------------
\begin{document}

% PORTADA ----------------------------------------------------------------------------------------------------------------------
\thispagestyle{empty}

\hspace{4.7cm}
\begin{minipage}{0.23\textwidth}
%\hspace{0.4cm}\crule[magenta]{0.5cm}{0.5cm}\hspace{0.1cm}\crule[magenta]{0.5cm}{0.5cm} 

%\vspace{0.01cm}

%\crule[magenta]{3.3cm}{0.35cm}

%\vspace{0.01cm}

%\crule[magenta]{3.3cm}{0.35cm} 
\end{minipage}
%---------------------------


%%%%%%%%%%%%%%%%%%%%% 
\vspace{1cm}
%%%%%%%%%%%%%%%%%%%%%%
%\begin{minipage}{0.38\textwidth}
%\noindent\textcolor{black}{\scalebox{1.35}[2.8]{ \bf Universidade da Coruña}}
%\end{minipage}
%------------------------------------
\hspace{3.4cm}
\begin{minipage}{0.25\textwidth}
%\centering
\includegraphics[scale=0.3]{Imagenes/03_Simbolo_logo_cor.png}
\end{minipage}

\large

\vspace{3cm}

\begin{center}
\begin{tabular}{rl}
 \colorbox{lightgray}{TÍTULO}     &	\titulouno         \\[10pt]
                                  &	\titulodos         \\[10pt]
                                  &	\titulotres        \\
                                  &                    \\
						 	                    &                    \\
							                    &                    \\
\colorbox{lightgray}{MÁSTER}			& \gradomaster       \\
% borrar la parte que no proceda
                                  &                    \\
                                  &                    \\
\colorbox{lightgray}{ASIGNATURA}  & \asignatura        \\
                                  &                    \\
                                  &  \\
							                    &  \\
							                    &  \\
\colorbox{lightgray}{ALUMNO}      &	\alumnouno         \\[10pt]
\colorbox{lightgray}{DIRECTOR}      &	\directoruno  \\       
									 & \directordos	\\	
                                  &  \\[10pt]


\colorbox{lightgray}{FECHA}       &	\fecha
\end{tabular}
\end{center}
\normalsize
% FIN DE LA PORTADA -----------------------------------------------------------------------------------------------------------------------------------
\newpage
% Índice -----------------------------------------------------------------------------------------------------------------------------------
\thispagestyle{plain}

% Índice de contenidos
\renewcommand{\contentsname}{Índice}
{\hypersetup{hidelinks}\tableofcontents}


% Listado de figuras
\newpage
\phantomsection
\renewcommand*\listfigurename{Listado de figuras}
{\hypersetup{hidelinks}\listoffigures}
\addtocontents{lof}{\protect\thispagestyle{fancy}}

% Listado de tablas
\phantomsection
\renewcommand*\listtablename{Listado de tablas}
{\hypersetup{hidelinks}\listoftables}
\addtocontents{lot}{\protect\thispagestyle{fancy}}

% Final del índice -----------------------------------------------------------------------------------------------------------------------------------

% Zona de resumenes en todos los idiomas -----------------------------------------------------------------------------------------------------------------------------------


\newpage


\section*{Resumen}

Resumen del trabajo desarrollado.

\section*{Resumo}

Resumo en Galego.

\section*{Abstract}

Abstract in English.

% Desarrollo del trabajo -----------------------------------------------------------------------------------------------------------------------------------

\newpage

\section{Introducción}

Contextualización del TFM desarrollado, describir los antecedentes y estado del arte que apliquen al mismo.


\section{Objetivos}
Describir los objetivos del trabajo, así como el grado de consecución de los mismos.

\section{Materiales y métodos}
Descripción de las tecnologías, herramientas, técnicas y conjuntos de datos 

\section{Experimentos / desarrollo}

Descripción del trabajo, así como los experimentos realizados.

\subsection{Sección nivel 2} 

\subsubsection{Sección nivel 3}
Se podrán incluir subsecciones de hasta tres niveles de profundidad, ejemplo:

\begin{enumerate}[2.]
    \item Materiales y métodos..
    \begin{enumerate}[2.1]
        \item Datasets.
        \begin{enumerate}[2.1.1]
            \item Dataset 1.
        \end{enumerate}
        \item Algoritmos empleados.
        \begin{enumerate}[2.2.1]
            \item Redes neuronales artificiales
        \end{enumerate}
    \end{enumerate}
\end{enumerate}

\section{Resultados}
Presentación de los resultados obtenidos en forma de tablas, gráficos, figuras, etc. y análisis de los mismos en base a datos objetivos presentados.

\section{Discusión}
(Sólo si es de aplicación en base a la temática del trabajo).
Comparación de los resultados obtenidos con investigaciones previas, mostrando los avances científicos obtenidos en relación al estado actual de las técnicas.

\section{Conclusiones y trabajos futuros}

\bibliographystyle{plain}
\bibliography{Bibliography}

Revisar la sección \ref{comentarios_referencias} donde se detalla el formato de las referencias

\section {Documentación de partida}
Documento de asignación del TFM donde se refleja la descripción del trabajo a realizar y los objetivos planteados.

\newpage
\section{Ejemplos de uso de tablas, figuras y ecuaciones}
En las siguientes subsecciones se muestra el uso de tablas, figura y ecuaciones en latex.


\subsection{Figuras}

Con los siguientes ejempos se describe el uso y configuración de imágenes en este documento de latex:




\begin{Verbatim}
\begin{figure}[h!] % comienza la figura
\centering % sirve para centrar horizontalmente la imagen en la página
\includegraphics{Imagenes/Fichero.ext} % se llama al fichero que contiene la imagen
\label{fig:etiqueta} % etiqueta que se usa para referenciar la figura
\caption[Nombre de la imagen en el listado de figuras]
        {Contenido del pie de la imagen}
\end{figure} % acaba la figura
\end{Verbatim}

Por ejemplo, al introducir el fichero \texttt{Simpsons.jpg} se obtiene la figura \ref{fig:Simpsons} situada en la página \pageref{fig:Simpsons}.

\bigskip

\begin{figure}[h!]
\hspace{1cm}\begin{minipage}{0.4\textwidth}
\begin{verbatim}
\begin{figure}[h!]
\centering
\includegraphics[scale=0.7]
    {Imagenes/Simpsons.jpg}
\caption[Familia Simpson]
        {Toda la familia junta}
\label{fig:Simpsons}
\end{figure}
\end{verbatim}
\end{minipage}
\begin{minipage}{0.5\textwidth}
\centering
\includegraphics[scale=0.7]
     {Imagenes/Simpsons.jpg}
\caption[Familia Simpson]
        {Toda la familia junta}
\label{fig:Simpsons}
\end{minipage}
\end{figure}

\bigskip

El comando \verb/[scale=0.7]/ permite el escaldo de la imagen si su tamaño no nos satisface. El valor del escalado es cualquier número (en nuestro caso $0.7$). Existe otra opción muy interesante, que permite recortar la figura. Esta es \verb/[trim= l b r t,clip]/ donde \verb/l/, \verb/b/, \verb/r/ y \verb/t/ especifican lo que hay que suprimir en cada lado de la imagen, siguiendo el orden izquierda, abajo, derecha y arriba (left, bottom, right, top).

\bigskip

Por ejemplo, si en la figura \ref{fig:Simpsons} se desea aislar a Burt Simpson se debe escribir

\begin{figure}[h!]
\centering
\begin{minipage}{0.6\textwidth}
\begin{verbatim}
\begin{figure}[h!]
\centering
\includegraphics[scale=0.7,trim= 165 20 50 55,clip]
    {Imagenes/Simpsons.jpg}
\label{fig:Burt}
\caption[Burt Simpson]
        {Burt Simpson}
\end{figure}
\end{verbatim}
\end{minipage}
\begin{minipage}{0.35\textwidth}
\centering
\includegraphics[scale=0.7,trim= 165 20 50 55,clip]
     {Imagenes/Simpsons.jpg}
\label{fig:Burt}
\caption[Burt Simpson]
        {Burt Simpson}
\end{minipage}
\end{figure}

\bigskip

También es posible incluir varias imágenes juntas, como vemos en la figura \ref{fig:miembros}, cuyo código es:
\begin{verbatim}
\begin{figure}[h!]
\centering
\subfigure[Marge]{\includegraphics[scale=0.09]{Imagenes/Marge.png}}
\hspace{0.3cm}
\subfigure[Burt]{\includegraphics[scale=0.35]{Imagenes/Burt.png}}
\hspace{0.3cm}
\subfigure[Homer]{\includegraphics[scale=0.10]{Imagenes/Homer.jpg}}
\hspace{0.3cm}
\subfigure[Lisa]{\includegraphics[scale=0.08]{Imagenes/Lisa.png}}
\label{fig:miembros}
\caption[Los miembros de la familia]{Los miembros de la familia} 
\end{figure}
\end{verbatim}

\bigskip

\begin{figure}[h!]
\centering
\subfigure[Marge]{\includegraphics[scale=0.09]{Imagenes/Marge.png}}
\hspace{0.3cm}
\subfigure[Burt]{\includegraphics[scale=0.35]{Imagenes/Burt.png}}
\hspace{0.3cm}
\subfigure[Homer]{\includegraphics[scale=0.10]{Imagenes/Homer.jpg}}
\hspace{0.3cm}
\subfigure[Lisa]{\includegraphics[scale=0.08]{Imagenes/Lisa.png}}
\caption[Los miembros de la familia]{Los miembros de la familia} 
\label{fig:miembros}
\end{figure}

\bigskip

Todas las figuras se numeran automáticamente con el formato:
\begin{center}
\textbf{Figura número\_de\_sección.número\_de\_figura}: Texto del pie de la figura
\end{center}


\subsection{Tablas}

A continución se muestra un ejemplo de uso de una tabla:

\begin{table}[h]
\begin{center}
\begin{tabular}{|c|r|l|}
\hline
Número  & Nombre    & Apellidos       \\ \hline
1       & Pedro & Picapiedra \\ \hline

\end{tabular}
\caption{Tabla de ejemplo}\label{ref_tabla}
\end{center}
\end{table}





Una tabla es un conjunto de celdas distribuidas en filas y columnas. Para su construcción, la herramienta básica es el entorno \texttt{tabular}, cuya estructura es:

\begin{center}
\begin{tabular}{cccccccccc}
\multicolumn{10}{l}{$\backslash$\texttt{begin}\{\texttt{tabular}\}\{\texttt{FormatoColumnas}\}} \\
\texttt{Celda 11} & \textcolor{red}{\texttt{\&}} &  \texttt{Celda 12} & \textcolor{red}{\texttt{\&}} & \texttt{Celda 13}  & \textcolor{red}{\texttt{\&}} & \texttt{...} & \textcolor{red}{\texttt{\&}} & \texttt{Celda 1p} &  \textcolor{blue}{\texttt{$\backslash\backslash$}} \\
\texttt{Celda 21} & \textcolor{red}{\texttt{\&}} &\texttt{Celda 22} & \textcolor{red}{\texttt{\&}} & \texttt{Celda 23} & \textcolor{red}{\texttt{\&}} & \texttt{...} & \textcolor{red}{\texttt{\&}} & \texttt{Celda 2p} &   \textcolor{blue}{\texttt{$\backslash\backslash$}} \\
\multicolumn{1}{c}{$\vdots$} &  &  \multicolumn{1}{c}{$\vdots$} &  & \multicolumn{1}{c}{$\vdots$} &  & \multicolumn{1}{c}{$\ddots$} &   & \multicolumn{1}{c}{$\vdots$}  &    \\
\texttt{Celda n1} & \textcolor{red}{\texttt{\&}} & \texttt{Celda n2} & \textcolor{red}{\texttt{\&}} & \texttt{Celda n3} & \textcolor{red}{\texttt{\&}} & \texttt{...} & \textcolor{red}{\texttt{\&}} & \texttt{Celda np} &   \textcolor{blue}{\texttt{$\backslash\backslash$}}  \\
\multicolumn{10}{l}{$\backslash$\texttt{end}\{\texttt{tabular}\}} \\
\end{tabular}
\end{center}



\bigskip 

\hangindent=1.3cm\hangafter=1 \textcolor{red}{\texttt{\&}}~ Marca el final de una columna y el principio de la siguiente. 

\medskip

\hangindent=1.3cm\hangafter=1 \textcolor{blue}{\texttt{$\backslash\backslash$}}~Es la instrucción utilizada para comenzar una nueva fila. 

\bigskip

\noindent El conjunto de códigos incluidos en \texttt{FormatoColumnas} indica el alineamiento de cada columna (especificadores) y la separación entre columnas (separadores).

\vspace{1cm}

\noindent\textbf{Especificadores}: son obligatorios y debe haber uno por cada columna. Estos son:

\begin{table}[h!]
\begin{tabular}{|lp{13cm}|}
\hline
\texttt{l}, \texttt{r}, \texttt{c} 
& Introduce una nueva columna justificada a izquierda, derecha o centrada respectivamente. La anchura de columna se determina automáticamente a partir de su contenido.
\\
\verb/p{Ancho}/
& Se utiliza en caso de que tengamos un texto largo en una de las columnas, creando una columna de anchura fija. Alinea por la parte superior del texto.
\\
\verb/m{Ancho}/
& Es similar al anterior, pero ahora la alineación se produce en la parte central del texto.\\
\hline
\end{tabular}
\caption[Especificadores]{Especificadores}
\end{table}

\vspace{1cm}

\noindent\textbf{Separadores}: son optativos, e indican cómo deben separarse las columnas.

\begin{table}[h!]
\begin{tabular}{|cp{13cm}|}\hline
 $|$ & Produce una barra vertical. \\
\verb/@{Objeto}/ & Suprime el espacio entre columnas e inserta en su lugar el \texttt{Objeto} declarado.\\
\hline
\end{tabular}
\caption[Separadores]{Separadores}
\end{table}


\bigskip

El primer ejemplo (tabla \ref{TS}) está formado tres columnas alineadas respectivamente a izquierda, centro y derecha. Los ejemplos que la siguen (tablas \ref{TP} y \ref{TM}) muestran una forma elegante de introducir párrafos en una tabla. En ellos se aprecia claramente la diferencia del funcionamiento de los códigos \verb/p{Ancho}/ y \verb/m{Ancho}/ para una columna concreta.



\begin{table}[h!]
\begin{minipage}{0.6\textwidth}
\begin{verbatim}
\begin{table}[h!]
\begin{tabular}{|l|c|r|}
\hline
izda & centro & dcha \\ \hline
Celda uno & Celda dos & Celda tres \\ \hline
\end{tabular}
\caption[Tabla simple]{Tabla simple}\label{TS}
\end{table}
\end{verbatim}
\end{minipage}
\begin{minipage}{0.4\textwidth}
\begin{tabular}{|l|c|r|}
\hline
izda & centro & dcha \\ \hline
Celda uno & Celda dos & Celda tres \\ \hline
\end{tabular}
\caption[Tabla simple]{Tabla simple}\label{TS}
\end{minipage}
\end{table}

\bigskip


\begin{verbatim}
\begin{table}[!ht]
\begin{tabular}{lp{14cm}}
\hline
\textbf{Mini} & Desde sus inicios en ... una marca.                     \\
              &                                                         \\
\textbf{Honda Civic} & La novena generación del ... audaz en el diseño. \\
\hline
\end{tabular}
\caption[Delimitador p\{~\}]{Delimitador p\{~\}}\label{TP}
\end{table}
\end{verbatim}

\bigskip

\begin{table}[!ht]
\begin{tabular}{lp{12cm}}
\hline
\textbf{Mini} & Desde sus inicios en 2001 Mini ha pasado de ofrecer un solo modelo a toda una amplia gama de interpretaciones diferentes que parece no tener final. Tanto es así que el nombre de Mini ha cambiado de ser un solo modelo a formar toda una marca.  \\
& \\
\textbf{Honda Civic} & La novena generación del Honda Civic se estrenó a comienzos de 2012. Algunos de los detalles de diseño futuristas se han marchado, pero a cambio cuenta con una calidad interior muy mejorada. En esta ocasión Honda ha sido menos audaz en el diseño. \\
\hline
\end{tabular}
\caption[Delimitador p\{~\}]{Delimitador p\{~\}}\label{TP}
\end{table}

\bigskip

\begin{verbatim}
\begin{table}[!ht]
\begin{tabular}{lm{14cm}}
\hline
\textbf{Mini} & Desde sus inicios en ... una marca.  \\
& \\
\textbf{Honda Civic} & La novena generación del ... audaz en el diseño. \\
\hline
\end{tabular}
\caption[Delimitador m\{~\}]{Delimitador m\{~\}}\label{TM}
\end{table}
\end{verbatim}

\bigskip

\begin{table}[!ht]
\begin{tabular}{lm{12cm}}
\hline
\textbf{Mini} & Desde sus inicios en 2001 Mini ha pasado de ofrecer un solo modelo a toda una amplia gama de interpretaciones diferentes que parece no tener final. Tanto es así que el nombre de Mini ha cambiado de ser un solo modelo a formar toda una marca.  \\
& \\
\textbf{Honda Civic} & La novena generación del Honda Civic se estrenó a comienzos de 2012. Algunos de los detalles de diseño futuristas se han marchado, pero a cambio cuenta con una calidad interior muy mejorada. En esta ocasión Honda ha sido menos audaz en el diseño. \\
\hline
\end{tabular}
\caption[Delimitador m\{~\}]{Delimitador m\{~\}}\label{TM}
\end{table}


\bigskip

La anchura de cada columna viene determinada por la longitud del mayor elemento incluido en las celdas de esa columna. 

\bigskip

En los ejemplos anteriores observamos que algunas filas y columna están separadas por líneas. Para insertar una línea horizontal por encima o por debajo de todas las celdas de una misma fila se usa el comando \verb/\hline/. Si queremos que la línea solo abarque algunas columnas, el comando adecuado es \verb/\cline{p-u}/, siendo \texttt{p} el número de la primera columna y \texttt{u} el de la última. Este comando se puede usar varias veces en la misma línea. Las líneas verticales que separan las columnas deben indicarse con el \texttt{Formato Columnas}.


\bigskip

A continuación introducimos algunos ejemplos que pueden ser de utilidad.



\begin{table}[!ht]
\centering
\begin{minipage}{0.6\textwidth}
\begin{verbatim}
\begin{table}[h!]
\begin{tabular}{|c|c|} \hline
\rowcolor{cyan} uno & dos \\ \hline
\rowcolor{green} tres & cuatro \\ \hline
\end{tabular}
\caption[Filas en color]{Filas en color}\label{FC}
\end{table}
\end{verbatim}
\end{minipage}
\begin{minipage}{0.3\textwidth}
\centering
\begin{tabular}{|c|c|} \hline
\rowcolor{cyan} uno & dos \\ \hline
\rowcolor{green} tres & cuatro \\ \hline
\end{tabular}
\caption[Filas en color]{Filas en color}\label{FC}
\end{minipage}
\end{table}

\medskip

\begin{table}[!ht]
\centering
\begin{minipage}{0.7\textwidth}
\begin{verbatim}
\begin{table}[h!]
\begin{tabular}
{|>{\columncolor{blue!25}}l|>{\columncolor{blue!10}}l|}
\hline
1 & 2  \\ \hline
3 & 4  \\ \hline
\end{tabular}
\caption[Columnas en color]
{Columnas en color}\label{CC}
\end{table}
\end{verbatim}
\end{minipage}
\begin{minipage}{0.2\textwidth}
\centering
\begin{tabular}
{|>{\columncolor{blue!25}}l|>{\columncolor{blue!10}}l|}
\hline
1 & 2  \\ \hline
3 & 4 \\ \hline
\end{tabular}
\caption[Columnas en color]
{Columnas en color}\label{CC}
\end{minipage}
\end{table}

\bigskip

Es posible agrupar varias celdas situadas en una misma fila, lo cual reducirá el número de columnas. El agrupamiento se realizará para cada fila de forma independiente. Se debe indicar en la orden de comienzo de la tabla el número máximo de columnas que habrá en la tabla. El comando que agrupa celdas es \verb/\multicolumn{número de columnas}{alineación}{Texto}/. 

\bigskip

\begin{center}
\begin{table}[!ht]
\begin{minipage}{0.65\textwidth}
\begin{verbatim}
\begin{table}[!ht]
\begin{tabular}{|c|c|c|c|c|c|c|} \hline
\multicolumn{7}{|c|}{Números primos} \\ \hline
\multicolumn{5}{|c|}{De 1 a 10} 
& \multicolumn{2}{|c|}{De 11 a 15} \\ \hline
1 & 2 & 3 & 5 & 7 & 11 & 13 \\ \hline
\end{tabular}
\caption[Comando multicolumn]{Comando multicolumn}
\label{TMC}
\end{table}
\end{verbatim}
\end{minipage}
\begin{minipage}{0.35\textwidth}
\begin{center}
\begin{tabular}{|c|c|c|c|c|c|c|}
\hline
\multicolumn{7}{|c|}{Números primos} \\ \hline
\multicolumn{5}{|c|}{De 1 a 10} & \multicolumn{2}{|c|}{De 11 a 15}\\
\hline
1 & 2 & 3 & 5 & 7 & 11 & 13 \\ \hline
\end{tabular}
\end{center}
\caption[Comando multicolumn]{Comando multicolumn}\label{TMC}
\end{minipage}
\end{table}
\end{center}

\subsection{Ecuaciones}

En la ecuación \ref{Equation_01}, se muestra un ejemplo de uso de ecuaciones.

\begin{equation}
e=\sum_{k=1}^{C}\sum_{x\epsilon Q_{k}}\left \|x-c_{k} \right \|^{2}
\label{Equation_01}
\end{equation}

\subsubsection{Ecuaciones matemáticas cortas}%%%%%%%%%%%%%%%%%%%%%%%%%%%%%%%%%%%%%%

Para indicar que se va a escribir una ecuación matemática se usa el símbolo \$, que se sitúa antes y después de la ecuación. 

\begin{center}
\texttt{La función \$\texttt{f(x)=sen(x)}\$ es periódica.} \hspace{1cm}
La función $f(x)=sen(x)$  es periódica.      
\end{center}  

\bigskip 

El texto que forma parte de la ecuación se escribe con estilo itálica. Para resaltar una ecuación frente al resto del texto se puede escribir aislada y centrada en una línea. El comando que nos permite hacer eso es: \verb/\begin{equation} EXPRESIÓN \end{equation}/.

\bigskip

\begin{center}
\begin{tabular}{m{8.5cm}m{7cm}}
\begin{verbatim}
El número $e$, cuyo valor decimal es:
\begin{equation}
e=2.718281
\end{equation}
no es un número racional. 
\end{verbatim}
&
El número $e$, cuyo valor decimal es:
\begin{equation}
e=2.718281
\end{equation}
no es un número racional.            
\end{tabular}
\end{center}

Observemos que automáticamente se asocia un número a la ecuación. Si lo que deseamos en centrar una expresión para resaltarla frente al texto, pero no queremos numerarla, lo que se debe escribir es:

\begin{center}
\begin{tabular}{m{8.5cm}m{6cm}}
\begin{verbatim}
El número $e$, cuyo valor decimal es:
$$ e=2.718281 $$
no es un número racional. 
\end{verbatim}
&
El número $e$, cuyo valor decimal es:
$$ e=2.718281 $$
no es un número racional.            
\end{tabular}
\end{center}

\bigskip

\subsubsection{Ecuaciones con más de una línea}%%%%%%%%%%%%%%%%%%%%%%%%%%%%%%%%%%%%%

Para escribir una ecuación matemática que ocupe más de una línea se usa el código 

\begin{verbatim}
                                  \begin{eqnarray}
                                  Líneas con ecuaciones
                                  \end{eqnarray}
\end{verbatim}

\noindent Estos códigos agrupan las fórmulas en tres columnas, separadas por el delimitador \&. Para cambiar de línea se usa \textbackslash\textbackslash. Todas las columnas se alinean por la columna central. Por ejemplo:


\begin{table}[ht!]
\begin{minipage}{0.5\textwidth}
\begin{verbatim}
\begin{eqnarray}
(a+b)(c+d) & = & a(c+d)+b(c+d) \\ 
           & = & ac+ad+bc+bd    
\end{eqnarray}
\end{verbatim}
\end{minipage}
%
\begin{minipage}{0.5\textwidth}
\begin{eqnarray}
(a+b)(c+d) & = & a(c+d)+b(c+d) \\ 
           & = & ac+ad+bc+bd    
\end{eqnarray}
\caption[Comando \texttt{eqnarray}]{Entorno \texttt{eqnarray}}
%
\end{minipage}
\end{table}

\bigskip

Aquí también vemos que se numeran la líneas automáticamente. Si deseamos no numerarlas se debe escribir \texttt{eqnarray*}. Si lo que deseamos es no numerar algunas, bastará con escribir \verb/\nonumber/ al principio de estas líneas.

\begin{table}[ht!]
\begin{minipage}{0.5\textwidth}
\begin{verbatim}
\begin{eqnarray}
\nonumber (a+b)(c+d) & = & a(c+d)+b(c+d) \\ 
                     & = & ac+ad+bc+bd    
\end{eqnarray}
\end{verbatim}
\end{minipage}
%
\begin{minipage}{0.5\textwidth}
\begin{eqnarray}
\nonumber (a+b)(c+d) & = & a(c+d)+b(c+d) \\ 
                     & = & ac+ad+bc+bd    
\end{eqnarray}
\caption[Comando \texttt{nonumber}]{Código \texttt{nonumber}}
%
\end{minipage}
\end{table}

\bigskip

 Una característica importante de \LaTeX\ es que numera las ecuaciones según el orden en que están escritas, por lo que si se modifica el texto o se añaden ecuaciones el procesador se encarga de numerarlas nuevamente. 

\vspace{1cm}

\section{Comentarios sobre las referencias}
\label{comentarios_referencias}

Todos los trabajos deben citar las fuentes consultadas. Estas suelen ser libros, revistas, páginas web o manuales de funcionamiento. La forma en que se citará cada tipo de referencia es:

\noindent\hangindent0.7cm\hangafter=1\textbf{Libro}: \textsc{Apellidos, Inicial del nombre}; \textit{Título del libro}, Edición, Ciudad de edición, Editorial, (Año de publicación).

\noindent\hangindent0.7cm\hangafter=1\textbf{Revista} \textsc{Apellidos, Inicial del nombre}; \emph{Título del artículo}, Revista, (Año de publicación), Volumen, Páginas.

\noindent\hangindent0.7cm\hangafter=1\textbf{Página web} \textit{Título del documento}, Organización, [Fecha de la consulta]. Disponible en: \url{http://www.pagina.ext/}

\noindent\hangindent0.7cm\hangafter=1\textbf{Manual} \textit{Título del documento}, Organización, [Fecha de la consulta]. Disponible en: \url{http://www.pagina.ext/}

						
\bigskip

Las referencias citadas siempre deben seguir un orden. El que se recomienda es el que afecta a las categorías mencionadas en la lista anterior. Dentro de cada categoría se ordenaran los libros y revistas atendiendo a la primera letra del primer apellido del/de la primer/a autor/a, las páginas web atendiendo a la primera letra del nombre de la página y los manuales atendiendo a la primera letra del nombre del manual \cite{CalvoRolle201537}. 



\end{document}


\end{document}